\begin{problem}[Compute sections on a principal open]\label{prob:vi17-p02}
Prove that for $X=\Spec A$ and $f\in A$, the canonical map $A_f\to\OO_X(D(f))$ is an isomorphism.
\end{problem}

\ifdefined\FullProblemDossiers
\paragraph{Why this problem matters.}
This is the computational heart of the affine structure sheaf.

\paragraph{Useful ideas.}
Construct the map by sending a fraction to its germs.  Injectivity is localization algebra.  For surjectivity, use a finite principal refinement of the local fraction cover together with the standard localization gluing argument.

\paragraph{Guided hints.}
\begin{enumerate}
\item A fraction $a/f^n$ defines a germ in every $A_{\mfp}$ for $\mfp\in D(f)$.
\item If its section is zero, use quasi-compactness of $D(f)$ or the localization zero criterion to obtain one power of $f$ annihilating $a$.
\item Cover a locally represented section by principal opens $D(g_i)\subseteq D(f)$.
\item Use compatibility in the overlaps to glue the localized elements.
\end{enumerate}

% VI17_FULL_SOLUTION_REFINED
% VI17_FULL_SOLUTION_REFINED
\begin{solution}
Define
\[
\Phi:A_f\longrightarrow\OO_X(D(f))
\]
by sending $a/f^n$ to the family of its images in the local rings $A_{\mfp}$ for $\mfp\in D(f)$.  Since $f\notin\mfp$ on $D(f)$, every power of $f$ is invertible in $A_{\mfp}$, so the family is well defined.  It is locally represented by the same fraction, hence is a regular section.  The construction is a ring homomorphism.

To prove injectivity, suppose $a/f^n$ maps to the zero section.  For each $\mfp\in D(f)$, the localization zero criterion gives $s_{\mfp}\notin\mfp$ with $s_{\mfp}a=0$.  The principal opens $D(s_{\mfp})$ cover $D(f)$.  Since $D(f)$ is quasi-compact, choose finitely many $s_1,\dots,s_r$.  The inclusion
\[
D(f)\subseteq D(s_1)\cup\cdots\cup D(s_r)
\]
implies $f^N\in(s_1,\dots,s_r)$ for some $N$.  Because every $s_i$ annihilates $a$, we obtain $f^Na=0$.  Hence $a/f^n=0$ in $A_f$.

For surjectivity, let $s\in\OO_X(D(f))$.  By local representability, $D(f)$ has an open cover on which $s$ is represented by fractions.  Refine it by principal opens $D(g_i)\subseteq D(f)$ and use quasi-compactness to choose a finite refinement.  On each $D(g_i)$ the section corresponds to an element of
\[
\OO_X(D(g_i))\cong A_{g_i}\cong (A_f)_{g_i/1}.
\]
On pairwise overlaps these localized elements agree because they are restrictions of the same section.  The standard finite localization equalizer for the ring $A_f$ therefore glues them to a unique element $u\in A_f$.  By construction $\Phi(u)$ and $s$ agree on the principal cover, hence are equal by the sheaf locality axiom.  Thus $\Phi$ is surjective and therefore an isomorphism.
\end{solution}



\paragraph{What happened mathematically.}
The topology and algebra are tuned to each other: principal opens are exactly the regions on which one chosen element has become invertible.

\paragraph{Extensions.}
Use the same argument for an $A$-module $M$ to prove $\widetilde M(D(f))\cong M_f$.

\paragraph{Related problems.}
VI.17.P1 and VI.17.P3 develop neighboring aspects of the same localization-sheaf calculus.

\paragraph{Provenance.}
\textbf{New canonical proof of the core VI/17 theorem, derived from the regular-function material in AG3 and CA13.}
\else
\paragraph{Provenance.} \textbf{New canonical proof of the core VI/17 theorem, derived from the regular-function material in AG3 and CA13.}
\fi
